% !TEX program = pdflatex
\documentclass[12pt]{article}
\usepackage[a4paper,margin=1in]{geometry}
\usepackage[T1]{fontenc}
\usepackage[utf8]{inputenc}
\usepackage{lmodern}
\usepackage{amsmath,amssymb,amsthm,mathtools,bm}
\usepackage{booktabs,array,tabularx}
\usepackage{enumitem}
\newcolumntype{L}[1]{>{\raggedright\arraybackslash}p{#1}}
\newcolumntype{Y}{>{\raggedright\arraybackslash}X}
\usepackage{setspace}
\usepackage[round,authoryear]{natbib}
\usepackage[colorlinks=true,linkcolor=blue,citecolor=blue,urlcolor=blue]{hyperref}

\onehalfspacing
\setlength{\parindent}{1.5em}
\setlength{\parskip}{0.12em}
\setlist[itemize]{leftmargin=2em,itemsep=0.25em,topsep=0.35em}
\setlist[enumerate]{leftmargin=2em,itemsep=0.25em,topsep=0.35em}
\emergencystretch=3em

\newtheorem{proposition}{Proposition}
\newtheorem{corollary}{Corollary}
\newtheorem{remark}{Remark}

\newcommand{\E}{\mathbb{E}}
\newcommand{\1}{\mathbf{1}}
\newcommand{\dd}{\mathrm{d}}
\newcommand{\calR}{\mathcal{R}}

\title{MAH, Commercialization Frictions, and Original-Drug Innovation}
\author{}
\date{}

\begin{document}
\maketitle
\thispagestyle{plain}

\begin{abstract}
\noindent
This paper develops a partial-equilibrium mechanism model of how China's Marketing Authorization Holder (MAH) system can affect original-drug innovation. Firms choose original-drug R\&D before successful opportunities are realized. After a successful opportunity arrives, a firm chooses whether to commercialize internally, retain authorization while entrusting production to a qualified manufacturer, transfer or out-license the opportunity, or abandon it. MAH affects expected R\&D value by adding a regulated holder-producer separation route, lowering the route-specific policy and contractual friction of that retained external route, and changing the composite probability that a latent opportunity becomes an observed approval or launch. The baseline closes only the qualified contract-manufacturing service market, generating an endogenous service price $p_m^*(M)$ that can attenuate the reform's effect. The model does not treat MAH as a demand shock, a scientific-productivity shock, a removal of holder responsibility, or a universal reduction in all commercialization costs. It delivers comparative statics for route value, R\&D intensity, and observed realized original-drug outcomes, and it clarifies which data moments can discipline composite objects in a mechanism calibration.
\end{abstract}

\noindent\textit{Keywords:} MAH; original-drug innovation; commercialization frictions; entrusted production; route choice.\\
\noindent\textit{JEL codes:} D22; L65; O31; O38; I18.

\section{Introduction}

China's MAH system separates the right to hold marketing authorization from a strict requirement of in-house production. This institutional change is potentially important for original-drug innovation because many research-originating firms can generate valuable projects but lack complete manufacturing capacity. The central mechanism studied here is therefore not demand expansion or scientific productivity. It is an expected-commercialization-value channel: when a firm anticipates that a future successful project can be commercialized through an additional retained route, current R\&D becomes more valuable.

The mechanism builds on two established ideas. First, pharmaceutical R\&D responds to expected downstream rewards \citep{Schmookler1966,AcemogluLinn2004,Finkelstein2004,DuboisDeMouzonScottMortonSeabright2015,BudishRoinWilliams2015}. Second, an innovator's ability to appropriate returns depends on complementary assets and commercialization strategy \citep{Teece1986,AroraFosfuriGambardella2001,GansStern2003}. MAH fits this second logic because it changes how an innovator may access manufacturing capability while retaining the authorization asset.

The model is deliberately narrow. It is not a full dynamic structural model of the pharmaceutical industry. It does not solve product-market clearing, R\&D input-market clearing, free entry, exit, or an invariant firm distribution. The baseline closes only the market directly tied to the MAH mechanism: the market for qualified contract-manufacturing services. This single closure is important because MAH-induced demand for entrusted production can raise the service price $p_m^*(M)$, preventing the model from mechanically implying that MAH increases innovation for all firms.

The contribution is to organize the mechanism in a way that is internally consistent with feasible data. The theory separates latent opportunity arrival from observed approvals or launches, keeps residual holder responsibility in the entrusted-route payoff, and treats route-use data as a proxy for realized holder-producer separation rather than as the policy treatment itself. The empirical role of the model is correspondingly modest: it maps feasible approval-side moments to composite objects such as route-realization value and innovation-realization scale, without claiming to identify every primitive friction separately.

The model is presented in two layers. The main text gives the shortest complete version needed for the paper's argument. The technical appendix gives the object-status audit, payoff accounting, Bellman derivation, CMO market regularity conditions, comparative-static derivations, and calibration boundaries. Research notes are kept outside the manuscript and are used only as a model-design control panel.

\section{Institutional Background}

The relevant institutional point is timing. Under China's drug-registration framework, a firm need not already be the final commercial holder when it starts R\&D. The economic choice is forward-looking: a firm invests today while anticipating the commercialization environment it will face if a project succeeds. After obtaining a Drug Approval License, the applicant becomes the MAH \citep{NMPARegistration2022}. The model therefore treats MAH as affecting expected future commercialization feasibility rather than current scientific productivity.

The MAH reform creates a regulated holder-producer separation route. The 2016 pilot allowed eligible applicants in selected regions to become holders under the pilot arrangement \citep{StateCouncilMAHPilot2016,CFDAMAHImplementation2016}. The later legal framework consolidated the principle that an MAH may manufacture by itself or entrust production to a qualified manufacturer \citep{DrugAdministrationLaw2019}. Subsequent NMPA rules require holder-side quality responsibility, assessment of the entrusted manufacturer, quality and entrustment agreements, and ongoing supervision \citep{NMPAHolderResponsibility2022,NMPAManufacturing2022,NMPAEntrustedProduction2023}.

These rules imply two modeling restrictions. First, entrusted production is not anonymous outsourcing. It is a regulated retained route in which authorization remains with the holder while production is performed by a qualified external manufacturer. Second, MAH does not eliminate holder responsibility. The holder retains quality, safety, lifecycle, and monitoring obligations. The model captures this through a residual entrusted-route burden $\mu_i^E$ and a composite realization probability $\zeta_i^E(M,p_m)$.

\section{Model}

\subsection{Environment and Timing}

There are two periods within the firm problem. In period $t$, firm $i$ chooses original-drug R\&D effort $x_i$. In period $t+1$, successful original-drug opportunities arrive, and the firm chooses a commercialization route for each opportunity. The institutional regime is $M\in\{0,1\}$, where $M=1$ denotes the MAH regime.

Firm characteristics are summarized by
\[
    \theta_i=(a_i,k_i,h_i^I,q_i^E,\mu_i^E,S_i,\tau_i^T),
\]
where $a_i$ is research productivity, $k_i$ indexes internal manufacturing capability, $h_i^I$ indicates whether internal production is feasible, $q_i^E$ indicates whether the opportunity is feasible for entrusted production, $\mu_i^E$ is the residual holder-side burden under entrusted production, $S_i$ is the gross non-retained transfer or out-licensing value, and $\tau_i^T$ is the transfer-route friction.

The model is partial equilibrium except for the qualified contract-manufacturing service market. It takes the market-return environment and firm-characteristic distributions as given. The service price $p_m^*(M)$ is determined by market clearing:
\begin{equation}
    D_m(p_m;M)=S_m(p_m).
    \label{eq:baseline_cm_clearing}
\end{equation}

\subsection{Route Set and Route Payoffs}

Conditional on a successful opportunity, the route set is
\begin{equation}
    \calR_i(M)=
    \{A,T\}\cup \{I:h_i^I=1\}\cup \{E:M=1,\ q_i^E=1\}.
    \label{eq:routeset}
\end{equation}
Here $I$ is internal production, $E$ is retained entrusted production, $T$ is a non-retained transfer, sale, or out-licensing option, and $A$ is abandonment or delay.

Let $R_i^{event}$ denote the route-independent commercialization-event return from a realized original drug, and let $v$ be the continuation value of adding one retained product or certificate to the firm's stock. The entrusted-route payoff is
\begin{equation}
    G_i^E(M,p_m)
    =
    \zeta_i^E(M,p_m)\left[\bar R_i^E+v\right]
    -p_m-\tau_i^E(M)-\mu_i^E,
    \label{eq:GE}
\end{equation}
where $\tau_i^E(M)$ is the MAH-reducible policy and contractual friction, and
\[
    \zeta_i^E(M,p_m)=\zeta^E(\tau_i^E(M),q_i^E,\theta_i,p_m)
\]
is a composite probability linking latent opportunity arrival to observed approval, launch, or realized retained outcome. The baseline uses $\bar R_i^E=\E[R_i\mid L_i=1,E]$ as a compressed post-production payoff object.

The other route payoffs are
\begin{align}
    G_i^I &= R_i^{event}+v-C^I(k_i), \label{eq:GI}\\
    G_i^T &= S_i-\tau_i^T, \label{eq:GT}\\
    G_i^A &=0. \label{eq:GA}
\end{align}
The transfer value $S_i-\tau_i^T$ is kept exogenous in the baseline. A Grossman-Hart-Moore-style outside-option interpretation is useful, but it is not added as a main state variable \citep{GrossmanHart1986,HartMoore1990}.

\subsection{Route Choice and CMO Market Closure}

Route choice follows a logit model throughout:
\begin{equation}
    P_i(r\mid M,p_m)
    =
    \frac{\exp(G_i^r(M,p_m)/\sigma_r)}
    {\sum_{\ell\in\calR_i(M)}\exp(G_i^\ell(M,p_m)/\sigma_r)},
    \qquad \sigma_r>0.
    \label{eq:logit_prob}
\end{equation}
The inclusive value of a successful opportunity is
\begin{equation}
    \Gamma_i(M,p_m)
    =
    \sigma_r\log
    \sum_{r\in\calR_i(M)}
    \exp(G_i^r(M,p_m)/\sigma_r).
    \label{eq:Gamma}
\end{equation}
The deterministic maximum is the limiting case as $\sigma_r\rightarrow0$, not the baseline model.

Demand for qualified entrusted manufacturing services is generated only by realized retained entrusted production:
\begin{equation}
    D_m(p_m;M)
    =
    \int
    a_i x_i^*(M,p_m)
    P_i(E\mid M,p_m)
    \zeta_i^E(M,p_m)
    \dd H_i .
    \label{eq:cm_demand}
\end{equation}
Qualified producer supply is upward sloping,
\begin{equation}
    S_m(p_m)=\int s(p_m,z_j^C)\dd H_C(z_j^C),
    \qquad
    \frac{\partial s(p_m,z_j^C)}{\partial p_m}\ge0.
    \label{eq:cm_supply}
\end{equation}
The equilibrium service price $p_m^*(M)$ solves equations \eqref{eq:baseline_cm_clearing}--\eqref{eq:cm_supply}. No other market is closed.

\subsection{R\&D Choice and Observed Outcomes}

R\&D effort produces latent successful original-drug opportunities according to
\begin{equation}
    \lambda_i^{lat}=a_i x_i,
    \label{eq:latent_arrival}
\end{equation}
with cost
\begin{equation}
    C_i^R(x_i)=\frac{\kappa}{2}x_i^2,\qquad \kappa>0.
    \label{eq:rd_cost}
\end{equation}
Because each latent success has expected route value $\Gamma_i(M,p_m^*(M))$, the firm solves
\[
    \max_{x_i\ge0}
    \beta a_i x_i\Gamma_i(M,p_m^*(M))-\frac{\kappa}{2}x_i^2 .
\]
The interior solution is
\begin{equation}
    x_i^*(M)=\frac{\beta a_i}{\kappa}\Gamma_i(M,p_m^*(M)).
    \label{eq:xstar}
\end{equation}
Observed realized output differs from latent opportunity arrival:
\begin{equation}
    \lambda_i^{obs}(M)
    =
    a_i x_i^*(M)
    \sum_{r\in\calR_i(M)}
    P_i(r\mid M,p_m^*(M))\zeta_i^r(M,p_m^*(M)).
    \label{eq:lambda_obs}
\end{equation}
For the retained entrusted route,
\begin{equation}
    \lambda_i^{obs,E}(M)
    =
    a_i x_i^*(M)
    P_i(E\mid M,p_m^*(M))\zeta_i^E(M,p_m^*(M)).
    \label{eq:lambda_obs_E}
\end{equation}

\section{Theoretical Implications}

\begin{proposition}[Logit-weighted R\&D response]
\label{prop:rd_response}
For a firm with $E\in\calR_i(M)$, holding $p_m$, $\zeta_i^E$, and other route payoffs fixed,
\begin{equation}
    \frac{\partial x_i^*}{\partial \tau_i^E}
    =
    -\frac{\beta a_i}{\kappa}P_i(E\mid M,p_m)<0.
    \label{eq:dxdtau}
\end{equation}
The effect of a lower entrusted-route friction is larger when the firm assigns higher route probability to retained entrusted production.
\end{proposition}

\noindent The result follows because the derivative of the log-sum inclusive value with respect to a route payoff equals that route's logit probability. The model therefore predicts probability-weighted sorting rather than a universal response. Firms for which the entrusted route is infeasible, unattractive, or dominated by internal production respond little.

The technical appendix also gives the full derivative when $\zeta_i^E$ varies with the entrusted-route friction. The main-text expression isolates the direct route-friction channel by holding realization probability fixed.

\begin{proposition}[Observed entrusted-route decomposition]
\label{prop:obs_decomposition}
The change in observed retained entrusted output can be decomposed as
\begin{align}
    \dd \lambda_i^{obs,E}
    =&\;
    a_i P_i(E\mid M,p_m^*(M))\zeta_i^E(M,p_m^*(M))\dd x_i^* \notag\\
    &+
    a_i x_i^*(M)\zeta_i^E(M,p_m^*(M))\dd P_i(E\mid M,p_m^*(M)) \notag\\
    &+
    a_i x_i^*(M)P_i(E\mid M,p_m^*(M))\dd\zeta_i^E(M,p_m^*(M)).
    \label{eq:obs_decomposition}
\end{align}
\end{proposition}

\noindent Observed approvals or launches may rise because firms invest more in R\&D, because successful projects are more likely to use the entrusted route, or because the composite realization probability improves. Approval-side data do not separately identify these components without additional application, status, and project-level information.

\begin{proposition}[CMO price attenuation]
\label{prop:cmo_attenuation}
Suppose MAH lowers $\tau_i^E$ for a positive-measure set of firms and raises demand for entrusted production. If qualified CMO supply is upward sloping, the induced increase in $p_m^*(M)$ partially offsets the entrusted-route value gain.
\end{proposition}

\noindent The entrusted-route payoff depends on the net wedge $p_m^*(M)+\tau_i^E(M)+\mu_i^E$. A reduction in policy or contracting friction increases route value only to the extent that it is not fully absorbed by a higher equilibrium service price or by residual holder-side responsibilities. This is the main force preventing a mechanical ``MAH raises innovation for everyone'' conclusion.

\begin{corollary}[Aggregate implication]
For a fixed firm distribution, aggregate latent original-drug opportunities rise when the positive $\Gamma_i(M,p_m^*(M))$ response among treated-relevant firms is not offset by CMO price feedback and when the affected set has positive measure. Aggregate observed original-drug outcomes additionally require that route-use and realization-probability terms do not fall enough to offset the latent-opportunity gain. The result is a private-value and realized-output implication, not a welfare theorem.
\end{corollary}

\section{Empirical and Calibration Implications}

The model is intended to discipline a mechanism calibration, not to identify every primitive friction. The key empirical distinction is between latent opportunity arrival and observed realized outcomes. Approval data, holder-producer links, and original-drug counts are close to realized outcomes. They are not direct observations of latent ideas, project-level scientific success, or true route-specific legal and managerial costs.

\begin{table}[htbp]
\centering
\footnotesize
\caption{Model Objects and Feasible Data Moments}
\label{tab:model_data_mapping}
\begin{tabularx}{\textwidth}{L{0.28\textwidth}Y Y}
\toprule
Model object & Feasible moment & Claim boundary \\
\midrule
$P_i(E\mid M,p_m^*(M))\zeta_i^E(M,p_m^*(M))$ & Holder-producer split among realized original drugs & Disciplines a composite route-realization wedge, not a primitive $\tau_i^E$ or $\mu_i^E$ separately. \\
$a_i x_i^*(M)$ and realization scale & Realized original-drug counts & Disciplines innovation-realization scale, not latent idea arrival alone. \\
$p_m^*(M)$ or CMO scarcity & Qualified entrusted-producer participation or capacity proxy & Needed to interpret attenuation; direct CMO prices may be unavailable. \\
Entry-cost distribution $F_e$ & Firm-year research-oriented entry panel & Future or sensitivity module unless credible entry data are merged. \\
\bottomrule
\end{tabularx}
\end{table}

Three calibration boundaries follow. First, route-use moments should be interpreted as evidence about relative route attractiveness among realized products, not as a direct estimate of the MAH policy treatment. Second, original-drug counts discipline a composite scale combining R\&D effort, research productivity, route choice, and realization probabilities. Third, entry should not be a baseline calibrated outcome unless the data contain a credible firm-year entry panel and pre-reform capability measures.

The baseline empirical exercise should therefore focus on approval-side realized-product moments: realized original-drug counts, original-drug shares, holder-producer splits, and heterogeneity by pre-reform manufacturing capability where available. Application-to-approval conversion, true project-level success probabilities, CMO prices, and research-oriented entry are natural future modules, but they should not be treated as already identified by approval-side records alone.

The formal appendix documents the corresponding moment equations, route-share log-odds mapping, and entry-module boundary. This separation is intentional: the main paper states what the feasible moments can support, while the appendix records the algebra behind that claim.

\section{Conclusion}

The paper's main mechanism is simple: MAH changes the value of R\&D by changing the commercialization route set available after a successful original-drug opportunity arrives. The retained entrusted route matters most for firms that can generate projects but lack complete internal manufacturing capability, and its value depends on route feasibility, residual holder responsibility, composite realization probability, outside options, and the equilibrium price of qualified contract-manufacturing services. The model remains intentionally partial equilibrium. This restriction is useful because it keeps the theory aligned with the paper's empirical objective: explaining how a specific institutional change can alter expected route value, R\&D incentives, and observed realized original-drug outcomes without turning the project into a full industry equilibrium or welfare model.

\clearpage
\bibliographystyle{aer}
\bibliography{mah_route_indicator_friction_refs}

\end{document}
